% 海森伯群（综述）
% license CCBYSA3
% type Sum

本文根据 CC-BY-SA 协议转载翻译自维基百科\href{https://en.wikipedia.org/wiki/Heisenberg_group}{相关文章}

在数学中，海森堡群$H$（以维尔纳·海森堡命名）是由如下形式的$3 \times 3$ 上三角矩阵组成的群：
$$
\begin{pmatrix}
1 & a & c \\
0 & 1 & b \\
0 & 0 & 1
\end{pmatrix}~
$$
其运算为矩阵乘法。元素$a, b, c$可以取自任意含幺交换环，常见的选择是实数环（得到所谓的“连续海森堡群”）或整数环（得到“离散海森堡群”）。

连续海森堡群出现在对一维量子力学系统的描述中，尤其是在Stone–von Neumann 定理的语境下。更一般地，可以考虑与$n$维系统相关的海森堡群，最一般的情况则对应于任意的辛向量空间。
\subsection{三维情形}
在三维情形下，两个海森堡矩阵的乘积为：
$$
\begin{pmatrix}
1 & a & c \\
0 & 1 & b \\
0 & 0 & 1
\end{pmatrix}
\begin{pmatrix}
1 & a' & c' \\
0 & 1 & b' \\
0 & 0 & 1
\end{pmatrix}
=
\begin{pmatrix}
1 & a + a' & c + ab' + c' \\
0 & 1 & b + b' \\
0 & 0 & 1
\end{pmatrix}.~
$$
由其中的$ab'$项可以看出，该群是非阿贝尔群。

海森堡群的中性元是单位矩阵，其逆元由下式给出：
$$
\begin{pmatrix}
1 & a & c \\
0 & 1 & b \\
0 & 0 & 1
\end{pmatrix}^{-1}
=
\begin{pmatrix}
1 & -a & ab - c \\
0 & 1 & -b \\
0 & 0 & 1
\end{pmatrix}.~
$$
该群是二维仿射群$\mathrm{Aff}(2)$的一个子群：$\begin{pmatrix}1 & a & c \\0 & 1 & b \\0 & 0 & 1\end{pmatrix}$作用在$(\vec{x}, 1)$上时，对应的仿射变换为：$\begin{pmatrix}1 & a \\0 & 1\end{pmatrix} \vec{x} +\begin{pmatrix}c \\b\end{pmatrix}$.三维情形下有若干重要的例子。
\subsubsection{连续海森堡群}
若$a, b, c$为实数（在实数环$\mathbb{R}$中），则得到连续海森堡群$H_{3}(\mathbb{R})$。

它是一个3维幂零实李群。

除了作为实$3 \times 3$矩阵的表示外，连续海森堡群在函数空间中还有若干不同的表示。根据Stone–von Neumann 定理，在同构意义下，海森堡群$H$存在唯一的不变的不可约酉表示，其中其中心通过一个给定的非平凡特征作用。

这一表示有若干重要的实现方式（或模型）：在薛定谔模型中，海森堡群作用在平方可积函数空间上；在θ表示中，它作用在上半平面上的全纯函数空间上；其名称源自其与θ函数的联系。
